
\exercise
For each of the following subsets of $\FF{3}$, determine whether it is a
subspace of $\F{3}$.
\begin{subtheorem}*
\item
$\br[\big]{ (x_1, x_2, x_3) \in \F{3}: x_1 + 2 x_2 + 3 x_3 = 0 }$

\item
$\br[\big]{ (x_1, x_2, x_3) \in \F{3}: x_1 + 2 x_2 + 3 x_3 = 4 }$

\item
$\br[\big]{ (x_1, x_2, x_3) \in \F{3}: x_1 x_2 x_3 = 0 }$

\item
$\br[\big]{ (x_1, x_2, x_3) \in \F{3}: x_1 = 5 x_3 }$
\end{subtheorem}
\exercisesubtheoremend


\exercise
Verify all assertions about subspaces in Example \ref{subspaces}.


\exercise
Show that the set of differentiable real-valued functions $f$ on the interval $(-4, 4)$ such that $f'\!(-1) = 3f(2)$ is a subspace of
$\RR{(-4, 4)}$.


\pagebreak[1]

\exercise
Suppose $b \in \R{}$. Show that the set of continuous real-valued functions $f$ on the interval $[0, 1]$ such that $\int_0^1 f = b$ is a subspace of $\R{[0, 1]}$ if and only if $b = 0$.


\exercise
Is $\R{2}$ a subspace of the complex vector space $\C{2}$?


\exercise
\begin{SUBtheorem}
\item
Is $\br[\big]{(a, b, c) \in \R{3}: a^3 =b^3}$ a subspace of $\R{3}$?
\item
Is $\br[\big]{(a, b, c) \in \C{3}: a^3 =b^3}$ a subspace of $\C{3}$?
\end{SUBtheorem}


\exercise
Prove or give a counterexample: If $U$ is a nonempty subset of $\R{2}$ such that $U$ is closed
under addition and under taking additive inverses (meaning $-u \in U$
whenever $u \in U$), then $U$ is a subspace of~$\RR{2}$.


\exercise
Give an example of a nonempty subset $U$ of $\R{2}$ such that $U$ is closed
under scalar multiplication, but $U$ is not a subspace of~$\RR{2}$.


\exercise
A function $\fun f {\R{}} {\R{}}$ is called \emph{periodic} if there exists a positive number $p$ such that $f(x) = f(x + p)$ for all $x \in \R{}$. Is the set of periodic functions from $\R{}$ to $\R{}$ a subspace of $\R{\R{}}$? Explain.


\exercise
Suppose $V\1_1$ and $V\1_2$ are subspaces of $V\1$. Prove that the intersection $V\1_1 \cap V\1_2$ is a subspace of $V\1$.


\exercise
Prove that the intersection of every collection of subspaces of $V$ is a subspace of $V\1$.


\exercise
Prove that the union of two subspaces of $V$ is a subspace of $V$ if and only if
one of the subspaces is contained in the other. \label{218}


\exercise
Prove that the union of three subspaces of $V$ is a subspace of $V$ if and only if one of the subspaces contains the other two.\label{1threesub}
\exercisecomment{This exercise is surprisingly harder than Exercise \ref{218}, possibly because this exercise is not true if we replace $\F{}$ with a field containing only two elements.}


\exercise
Suppose
\[
U = \br[\big]{ (x, -x, 2x) \in \F3 : x\in \F{}} \andd W = \br[\big]{ (x, x, 2x) \in \F3: x \in \F{}}.
\]
Describe $U+W$ using symbols, and also give a description of $U + W$ that uses no symbols.


\exercise
Suppose $U$ is a subspace of $V\1$.  What is $U + U$?


\exercise
Is the operation of addition on the subspaces of $V$ commutative?
In other words, if $U$ and $W$ are subspaces of $V\1$, is
$U + W = W + U$?\index{commutativity}


\exercise
Is the operation of addition on the subspaces of $V$ associative?
In other words, if $V\1_1,V\1_2, V\1_3$ are subspaces of $V\1$, is
\[
(V\1_1 + V\1_2) + V\1_3 = V\1_1 + (V\1_2 + V\1_3)?
\]
\exercisedisplayend


\exercise
Does the operation of addition on the subspaces of $V$ have an
additive identity?  Which subspaces have additive inverses?


\exercise
Prove or give a counterexample: If $V\1_1, V\1_2, U$ are subspaces of~$V$ such
that
\[
V\1_1 + U = V\1_2 + U,
\]
then $V\1_1 = V\1_2$.


\exercise
Suppose
\[
U = \br[\big]{(x, x, y, y) \in \F{4}: x, y \in \F{}}.
\]
Find a subspace $W$ of $\F{4}$ such that
$\F{4} = U \oplus W\3$.


\exercise
Suppose
\[
U = \br[\big]{(x, y, x+y, x-y, 2x) \in \F{5}: x, y \in \F{}}.
\]
Find a subspace $W$ of $\F{5}$ such that
$\F{5} = U \oplus W\3$.


\exercise
Suppose
\[
U = \br[\big]{(x, y, x+y, x-y, 2x) \in \F{5}: x, y \in \F{}}.
\]
Find three subspaces $W\3_1, W\3_2, W\3_3$ of $\FF{5}$, none of which equals $\{0\}$, such that
$\F{5} = U \oplus W\3_1 \oplus W\3_2 \oplus W\3_3$.


\exercise
Prove or give a counterexample: If $V\1_1, V\1_2, U$ are subspaces of~$V$ such that
\[
V = V\1_1 \oplus U \andd V = V\1_2 \oplus U,
\]
then $V\1_1 = V\1_2$.
\hint{When trying to discover whether a conjecture in linear algebra is true or false, it is often useful to start by experimenting in\/ $\FF2$.}


\exercise
A function $f\colon \R{} \to \R{}$ is called \emph{even} if
\[
f(-x) = f(x)
\]
for all $x \in \R{}$. A function $f\colon \R{} \to \R{}$ is called \emph{odd} if
\[
f(-x) = -f(x)
\]
for all $x \in \R{}$. Let $V\1_\text{e}$ denote the set of real-valued even functions on $\R{}$ and let $V\1_\text{o}$ denote the set of real-valued odd functions on $\R{}$.
Show that $\R{\R{}} = V\1_\text{e} \oplus V\1_\text{o}$.


